\section{Introduction}
In our project, we investigate how a 16-bit RISK-V architecture works.
We explore how simple arithmetic operations like multiplication and division are performed, their synthesis in combinatorial logic, and how other parts of the processor interact with each other.
We have done work to adjust the default structure of the RISK-V processor so that it matches the given specifications.
This report will give the reader a better understanding and an overview of RISK-V architecture.
The source code to the Verilog implementation of it and its test benches can be found online\footnote{At \url{https://github.com/superstudy160/HDL-RISC-V-i16-A-Project/}.}.

\subsection{Notation and terminology}

\begin{itemize}
    \setlength\itemsep{3pt}
    \item $R_i$ -- the $i$-th register of the processor;
    \item \codeword{MonoFont} -- the indication of a code snippet, parameter name of a keyword;
    \item \codeword{l} and \codeword{lv} -- parameters used throughout our code to specify the bitness of the parameters. By default, \codeword{l} = 16 (given that we are working with a 16-bit processor), and \codeword{lv} = 15 (the highest bit index of a 16-bit number). But we can decrease them for the testing purposes.
\end{itemize}